\levelstay{Notation}
This chapter is about notation.

\begin{itemize}
\item Quantum operators have hats like this: $\ahat$.

\item Matrices (whether quantum or classical) are bold text like this: $\mathbf{C}$. Linear transformations are in normal math text.

\item Anharmonicity is $\eta$ and is always added (so transmons have $\eta < 0 $).

\item $\hat{H} / \hbar =$  expression with dimensions of angular frequency. Example: $\hat{H} / \hbar = \Omega \sz$.

\item Linear frequencies are quoted in Hz. Angular frequencies and rates are quoted in $\mathrm{s}^{-1}$. For instance, a transmon may have a qubit frequency of 4 GHz and a relaxation rate of $5\times10^3 \ \mathrm{s}^{-1}$.

\item We will usually write equations using quantities that are typically defined as angular frequencies, but quote quantities in linear frequencies units. So we will write a bare qubit Hamiltonian like $\hat{H} = -\frac{\omega_q}{2}\sz$, but then say the qubit frequency is 4 GHz. We will clarify this by explicitly including the factor of $2\pi$, as in $\omega_q / 2\pi = 4\,\text{GHz}$.

\item Maps and vector spaces will use calligraphic fonts: $\mathcal{H}$. 

\item Fourier transform: $f(t) = \int \frac{d\omega}{2 \pi} \tilde{f}(\omega) e^{i \omega t}$.

\item The imaginary number is $i$.

\item Words or abbreviations in subscript or superscript are not italicized like this: $\omega_\mathrm{diss}$
\end{itemize}